% ============================================================================
% sections/method.tex -- Section 6: Experimental Method
% Drafted from the P1-approved docs/prereg.md and docs/method.md (§7, §8, §9).
% DESIGN ONLY: this section contains no results and no measured numbers.
% ============================================================================

\section{Experimental Method}
\label{sec:method}

The evaluation is a staged campaign of scripted faults, verified by an independent
oracle and analysed under a pre-registered protocol. This section states the substrate,
the injectors, the oracle and its validation gates, the design of the stages, the
metrics, and the frozen statistical plan. All design constants (cell counts,
repetitions, quiescence bounds, test thresholds) were fixed before the first
measurement; the protocol and the analysis code are archived with the artifact.

\subsection{Implementation and substrate}
\label{sec:method:substrate}

The artifact is a Java~21 service (one process per consumer instance) that implements all
nine configurations behind two flags: the mechanism and the commit ordering
(\emph{strict} or \emph{early}). The load generator is a Java program that emits the
deterministic event stream and writes the ground truth; the oracle is an independent
Python program that shares no code with the service. The service emits structured phase
markers and writes its effective configuration at startup and its counters at shutdown.

The reported campaign runs on a \emph{native-process substrate}: three Kafka~3.7.2 KRaft
brokers, PostgreSQL~16.11, Toxiproxy~2.9.0 and the service, all as native Windows
processes on one host, with JDK~21.0.13.1 and Maven~3.9.9. This is an explicit deviation
from a containerised deployment: the host has no Docker, no administrator rights and no
WSL2. The substrate has three consequences that the paper carries as limitations.
First, isolation is process-level: CPU affinity, priority class, heap bounds, connection
pools and every broker/database setting are pinned and recorded per run, and the shared
configuration is hashed so that a non-intrinsic difference between mechanisms fails the
run. Second, the Kafka cluster is three processes on one operating system rather than
three machines: a broker kill is a real process death and a real leader election, but
not a machine loss. Third, container-level guarantees (cgroup limits, image digests) are
replaced by recorded process limits and SHA-256 digests of the vendored binaries, and
the periodic full-stack restart every 50 runs bounds cross-run interference. Each run
uses a fresh logical state --- new source and publication topics, a new consumer group,
and a fresh database schema --- and the identities of every process are recorded so that
an injector can target a specific one.

\subsection{Fault injection}
\label{sec:method:injection}

Two injection levels are used, because they answer different questions. \emph{Level~1}
places a fault deterministically at a named crash point of Section~\ref{sec:model}:
failpoints at fixed code locations trigger one of four actions (crash, throw, abort the
transaction, or pause) on a chosen occurrence. Level~1 is the structural experiment: a
mechanism that survives a fault placed at its worst location has absorbed it.
\emph{Level~2} applies the six scripted faults of Table~\ref{tab:faults} at random,
recorded times within the steady-state window; it measures how often a fault realises an
exposed window in a realistic timeline. Every armed injection writes two attestation
records to the fault log --- one when it fires and one when the system is restored ---
and a run whose injection did not fire is invalid, excluded and counted, with the
exclusion rate reported per cell.

\begin{table*}[t]
\centering
\small
\caption{Scripted faults (Level~2) and the crash classes they realise. Timing is either
phase-coupled (Level~1) or a seeded random offset (Level~2); every action is attested in
the fault log. F4 is split by commit direction; F6 is protocol-level duplicate
redelivery only.}
\label{tab:faults}
\begin{tabular}{@{}p{0.7cm}p{5.4cm}p{4.4cm}p{5.2cm}@{}}
\toprule
Fault & Injected condition & Mechanism manipulated & Crash classes \\
\midrule
F1 & one broker killed mid-produce (leader and transaction coordinator of the instance) & Kafka broker process, hard kill & $K_2$, $K_3$ \\
F2 & consumer instance killed during processing, forcing a rebalance & consumer process, hard kill & $K_1$--$K_3$, $K_5$ \\
F3 & produce or transaction commit times out with unknown outcome & TCP delay/timeout on the produce link & $K_2$, $K_3$ \\
F4a & effect commit durable, progress commit fails & offset-commit link cut after the effect marker & $K_3$ \\
F4b & progress commit durable, effect commit fails (early ordering probe) & ordering flag plus a database-link cut & $K_4$ \\
F5 & network partition between the service and the database & TCP blackhole/reset on the database link & $K_2$, $K_3$ \\
F6 & duplicate redelivery after a rebalance (protocol-level replay) & consumers stopped, resume point rewound, consumers resumed & $K_5$ \\
\bottomrule
\end{tabular}
\end{table*}

\subsection{Independent oracle and ground truth}
\label{sec:method:oracle}

Correctness is measured against ground truth created \emph{before} the effect: the
generator derives the whole event stream from the repetition seed, writes one canonical
line per expected posting to a ground-truth file, computes its SHA-256, and fsyncs both
before producing the first record. After the run, the oracle (i) regenerates the expected
posting set from the seed with its own implementation, (ii) checks that the regenerated
digest equals the pre-send digest, (iii) dumps the source topic and checks that it is
contiguous and exactly matches the produced manifest, and (iv) compares the sink against
the expected multiset. Only when the first three checks pass is $\Omega$ well defined
and the run eligible for analysis. The oracle reads the sink through a read-only role
and computes the invariants I1--I4 of Table~\ref{tab:invariants} and the verdict of
Eq.~\eqref{eq:consistency} directly from the dumps.

Two properties of the oracle protect the conclusions. The first is independence: it is a
separate program in a different language, with no shared code, and ground truth exists
before the pipeline runs. The second is mutation testing: before any measurement, the
checker is attacked with seeded synthetic defects --- a dropped posting, a duplicated
posting, a flipped leg sign, an updated row, and a spurious posting, at least one hundred
mutations per class --- and it must detect one hundred percent of them, attributed to the
correct invariant. A checker that misses a defect class invalidates the campaign rather
than a mechanism. The approach follows the systematic-testing tradition of crash
consistency \cite{crashmonkey2019} and isolation-anomaly checking \cite{elle2020}: the
verifier is an independent executable model of the property, not an assertion inside the
system under test.

\textbf{Quiescence.} A run is observed only in a quiescent state: consumer lag is zero
for all source partitions, no new sink row has appeared for ten seconds, no application
transaction is active, and the retry drain interval has elapsed. If quiescence is not
reached within a pre-registered bound of 300~s, the run is \emph{censored}: it counts as
a violating run (or as a quiescence failure if the application is still active) rather
than being dropped. The final observation is a single repeatable-read snapshot
\cite{pg-isolation}.

\subsection{Harness-validation gates}
\label{sec:method:gates}

No measurement is analysed until the gates of Table~\ref{tab:gates} pass. The gates are
deliberately adversarial: the negative control must fail, the positive controls must pass,
and the oracle must catch synthetic defects before it is trusted to judge the pipeline.
The gate that most often reveals a design error is the negative control: M0 is expected
to produce duplicates under the deterministic $W_1$ injection, so a clean M0 means the
injector or the oracle is broken, not that naive at-least-once is safe.

\begin{table*}[t]
\centering
\small
\caption{Harness-validation gates. G1--G8 must pass on pilot data before any experimental
run is analysed; failures halt the campaign.}
\label{tab:gates}
\begin{tabular}{@{}p{0.7cm}p{8.3cm}p{6.4cm}@{}}
\toprule
Gate & Requirement & Failure action \\
\midrule
G1 & negative control: M0 shows duplicates under the deterministic $W_1$ injection, with attestation & halt; harness broken \\
G2 & positive controls: M1-kafka and M4+M3 show no violation in fault-free runs & halt \\
G3 & checker mutation testing: 100\,\% detection of drop, duplicate, sign flip, update, spurious row ($\geq 100$ mutations per class) & fix checker; repeat \\
G4 & injector attestation: one fired and one restored record per armed injection; exclusion rate computed from receipts only & exclude and count; flag cell if $>20\,\%$ \\
G5 & ground-truth independence: generator digest equals oracle digest; source-integrity gate passes & invalidate runs; fix derivation \\
G6 & configuration equality: shared configuration hash identical across mechanisms; only declared intrinsic keys differ & halt; unfair comparison \\
G7 & load-generator capability: sustained $\geq 3\times$ the fastest measured mechanism throughput; doubling the quiet period does not change any count & resize generator; recalibrate quiescence \\
G8 & clock and host sanity: single boot, no clock step, free disk above the floor, affinity and priority as pinned & flag or rerun \\
\bottomrule
\end{tabular}
\end{table*}

\subsection{Design, metrics, and statistics}
\label{sec:method:design}

The campaign is staged so that the structural questions are answered cheaply and
deterministically before the frequency and performance questions are measured under load
(Table~\ref{tab:stages}). Stage~A runs every reachable crash class through every
mechanism with deterministic Level-1 injections at small load; its cells are the 52
reachable (mechanism, crash-class) pairs of Table~\ref{tab:crashverdicts} plus one
control cell per mechanism. Stage~B applies the six scripted faults at random times to
all mechanisms at the canonical operating point and adds the sensitivity blocks. Stage~C
measures performance in the fault-free case; Stage~D reuses Stage~B runs for performance
under faults and for time-to-consistent-state. The campaign was pre-registered with an
unattended wall-clock budget of 8--24 hours; if the pilot projection exceeds the budget,
a fixed cut ladder is applied in order --- the large-load block first, then the partition
sensitivity block, then repetitions down to the pre-registered floor for correctness
cells; control cells are never cut.

\begin{table*}[t]
\centering
\small
\caption{Staged design. Repetitions $N$ apply per cell; ``cells'' counts executed
(mechanism, condition) pairs, with structurally unreachable cells recorded as not
applicable rather than as passes. The total is 7{,}866 matrix runs plus 20 validation
runs. No number in this table is a measurement.}
\label{tab:stages}
\begin{tabular}{@{}p{0.7cm}p{4.6cm}p{5.0cm}p{1.4cm}p{1.5cm}@{}}
\toprule
Stage & Purpose & Design & $N$ & Runs \\
\midrule
P & calibrate per-run cost, validate oracle and injectors, measure capacity for the target-rate rule & 9 mechanisms $\times$ 8 conditions $\times$ 1k events, 3 partitions & 3 & 216 \\
A & structural correctness under deterministic crash points & 52 crash-class cells + 9 control cells $\times$ 1k events, 3 partitions & 30 & 1{,}830 \\
B & correctness under stochastic faults at realistic load, plus sensitivity & 67 fault cells $\times$ 100k events, 6 partitions; sensitivity blocks at 10k events & 30 / 10 & 4{,}020 \\
C & performance, fault-free & 9 mechanisms $\times$ 2 loads $\times$ 3 partition counts; large-load and capacity blocks & 30 / 10 & 1{,}800 \\
D & performance under faults and time-to-consistent-state & reuse of Stage~B runs & --- & 0 \\
\bottomrule
\end{tabular}
\end{table*}

\paragraph{Metrics.}
Correctness is measured per run by the counts of duplicate postings ($\mathrm{I3}$),
lost postings ($\mathrm{I4}$), imbalanced entries ($\mathrm{I1}$) and mutated postings
($\mathrm{I2}$), and by the binary outcome \emph{violating run} --- any of the four
counts positive. Performance is measured in two modes. In the \emph{fixed-rate} mode the
generator offers events at a target rate independent of the service, so latency is not
corrupted by coordinated omission; in the \emph{saturation} mode the offered rate is far
above capacity and the outcome is achieved throughput. End-to-end latency of an event is
the interval from its scheduled emission to the observation that the effect commit for
its batch returned; per-run percentiles p50, p90, p99 and p99.9 are recorded with
HdrHistogram \cite{hdrhistogram} at a pre-registered precision and range. The first
$\max(2000,\,10\,\%)$ events of a run are discarded from performance statistics as
warm-up and retained for correctness. Throughput is committed events over the
steady-state window. Time-to-consistent-state is the interval from the end of load to the
first quiescent sample, and is reported with the censoring rate. Measurement practice
follows the benchmarking literature: fixed machine state and identical configuration for
all mechanisms, median and IQR rather than single observations, tail percentiles rather
than means \cite{dean2013}, awareness that small configuration differences can reverse
conclusions \cite{mytkowicz2009,raasveldt2018}, explicit reporting of the measurement
setup \cite{hoefler2015}, and statistically rigorous treatment of JVM effects
\cite{georges2007}. The design adopts the benchmarking discipline advocated for stream
processing engines \cite{theodolite2021}.

\paragraph{Frozen statistics.}
The analysis plan was frozen, together with the analysis code, before the first
experimental run; anything not in the plan is labelled exploratory. For correctness,
each cell reports the proportion of violating runs with an exact
Clopper--Pearson 95\,\% confidence interval \cite{brown2001,jsta2014}; when no violation
is observed, the rule-of-three bound is reported as a conservative summary
\cite{rule3} and the text states explicitly that the bound does not prove absence.
Mechanisms are compared with McNemar's exact test on paired runs --- the same repetition
seeds are used for every mechanism within a cell, so the comparison is paired
\cite{mcnemar1947} --- with Fisher's exact test as an unpaired sensitivity analysis, and
Holm's correction within pre-declared families \cite{holland1987}. For performance, each
cell reports median and IQR with a bootstrap 95\,\% confidence interval (10{,}000
resamples, recorded seed) \cite{efron1979}; mechanisms are compared with the Friedman
test over blocks of load and partition count, Nemenyi post-hoc analysis, and Cliff's
delta with pre-registered magnitude thresholds (negligible below 0.147, small below 0.33,
medium below 0.474) \cite{demsar2006,cliff1993}. The level is $\alpha=0.05$ throughout,
unadjusted and adjusted $p$-values are both reported, all $N$ and exclusion rates are
reported per cell, and power statements accompany zero-event claims. Statistical
practice follows the guidance for randomised and benchmark comparisons in software
engineering \cite{arcuri2011,arcuri2012}, and the raw data, seeds and frozen scripts
are released so that every table can be regenerated \cite{collberg2016}.
